\section{Discussion and Limitations}
\label{sec:discussion}

\paragraph{When to use a one-step \nom{} world model.}
The results reduce to one criterion. If the outcome uncertainty consists of a few discrete alternatives, as in bouncing directions, a hidden centre of gravity or a branching obstacle, the random codebook can cover them: the distortion $I(m)$ vanishes exponentially once $m$ exceeds the number of modes on the outcome manifold, a budget of $16$--$32$ removes the memorization incentive, and the model covers the modes with one forward pass, stays on the token manifold under rollout and estimates equiprobable branch risks correctly. If instead the uncertainty is continuous and high-dimensional, as in the texture and fine motion of real video, the codebook cannot cover it at any affordable $m$; the samples collapse onto a few context-specific directions and the generator is best read as an accurate point predictor with a little aligned spread. In between, the budget distorts mode probabilities in a known way. A practical rule for choosing $m$ is to estimate $I(m)$ from the outcome diversity of the training data, for instance from a simulator that enumerates branches or from the outcomes of near-identical contexts, and to pick the smallest $m$ at which $I(m)$ is a few percent of the mode spacing. A hybrid that uses \nom{} to select the mode and a small stochastic head or a one- or two-step flow for the residual is the natural way to cover both regimes; we leave it to future work.

\paragraph{Limitations.}
(1) The approach under-disperses on real video and no variant in the one-step family fixed it; the quantitative agreement of the isotropic model on BAIR uses one fitted parameter, and the independent estimate of that parameter failed. (2) On stochastic Push-T the advantage over flow matching is an advantage in precision; the recall advantage does not clear seed variability with three seeds, and a shared $12$\,px contact error masks differences. (3) The benchmarks are synthetic or simulated except BAIR; the models are small, and our flow baselines, although trained with FlowWM's recipe and checked for under-training, are far from the scale of \citet{porcher2026flowwm}, so the comparison is confined to matched small-scale compute. (4) Planning success did not improve on either task, even where the model's risk estimates were accurate. (5) Large budgets flatten mode probabilities; the samples of such a model are a coverage device and need a separate calibration mechanism, which the energy score did not provide in non-repeating contexts. (6) The ablations, the BAIR variants and the deterministic Push-T control are single-seed; the $K{=}2,16$ grid and the energy-score follow-ups are two-seed. (7) Proposition~\ref{prop:es} was derived after we observed the energy score memorizing; only the follow-up runs (A2-fix, per-token score) had predictions recorded in advance, and we examined one $\beta$ and two distance choices, not mixtures of the two objectives. (8) The mathematics of Section~\ref{sec:theory} is classical; the contribution is the application and its tests.

\section{Conclusion}
\label{sec:conclusion}
A one-step world model trained with a \nom{} objective on frozen DINO tokens covers enumerable futures better than multi-step flow matching and other one-step generators at matched parameters, with one network evaluation per step and rollouts that stay on the data manifold. The objective is the distortion of a random codebook of $m$ draws, and that single fact answers the question in the title: the budget must make the codebook's miss distortion small relative to the spacing of the outcome modes, which a few dozen candidates achieve for mode-like uncertainty and no budget achieves for continuous high-dimensional uncertainty. The same accounting shows why, in contexts that never repeat, a strictly proper energy score memorizes the realized outcome while \nom{} can be made not to, why mode probabilities flatten with the budget, and where the method ends. We hope the memorization index and the enumerable benchmarks are useful to others building one-step stochastic world models, whichever objective they choose.
